% !TEX root = ../main.tex

%----------------------------------------------------------------------
\subsection{Why a compact model suffices}
%----------------------------------------------------------------------

\dart{}'s 0.75\,M parameters outperform ALIGNN's 4.03\,M by a substantial
margin.
This result aligns with emerging evidence that, in the data-limited regime
characteristic of materials science ($N \sim 10^3$--$10^4$), training
strategy and data engineering dominate over architectural complexity.
Our scaling analysis (not shown) finds a power-law exponent
$\alpha = -0.40$ for data scaling but $\beta \approx -0.01$ for parameter
scaling ($R^2 = 0.95$), confirming that \emph{more data, not more
parameters}, is the primary lever for improvement on this task.

%----------------------------------------------------------------------
\subsection{Self-attention as learned potential summation}
%----------------------------------------------------------------------

The success of the global geometric attention layers can be understood
through the lens of ``neural potential summation''~\cite{crystalformer2024}:
the learned distance bias $\phi_\mathrm{bias}(d_{ij}^{\mathrm{PBC}})$ acts
as a non-parametric interatomic potential that gates information flow.
For defect systems, where the perturbation source (dopant) induces
long-range elastic strain fields decaying as ${\sim}1/r^2$ in 2D, this
mechanism naturally allocates higher attention weights to nearby atoms while
retaining non-zero coupling to distant atoms---precisely the physics of
defect-induced perturbation.

The attention-weight analysis (\S\ref{sec:interpret}) provides empirical
confirmation: the defect atom emerges as a global attention hub without
explicit architectural enforcement, consistent with its role as the dominant
perturbation source.

%----------------------------------------------------------------------
\subsection{Methodological value of graduated OOD assessment}
%----------------------------------------------------------------------

Our three-tier protocol addresses a blind spot identified by
Li~\textit{et al.}~\cite{li2025ood}: most materials ML papers report only
random-split performance, which tests interpolation rather than genuine
extrapolation.
By systematically varying the degree of distributional shift---from
in-distribution to intra-family host transfer to compositional
extrapolation---we provide a nuanced picture of model reliability.

We recommend that future materials ML studies adopt similar tiered
protocols, with at least one OOD tier beyond random splitting.
The LOHO framework is readily transferable to other multi-component
material families (e.g., high-entropy alloys, hybrid perovskites).

%----------------------------------------------------------------------
\subsection{Limitations and outlook}
%----------------------------------------------------------------------

Several limitations warrant acknowledgment:
(1)~The current model covers only the 44 host materials in IMP2D;
generalization to entirely new material families (e.g., MXenes, h-BN
alloys) requires expanded training data.
(2)~We predict only defect formation energies; extending to other
defect-related properties---charge transition levels, migration barriers,
optical transition energies---would broaden the platform's utility for
defect engineering.
(3)~The OOD evaluation is limited to the G6-TMD subfamily; a broader
cross-family assessment awaits datasets with sufficient diversity.

Looking forward, the integration of \dart{} with active learning loops
could efficiently expand the training distribution into unexplored
regions of host--dopant space, guided by the calibrated uncertainty
estimates demonstrated here.
